\documentclass[11pt]{article}

\usepackage[utf8]{inputenc}
\usepackage[T1]{fontenc}
\usepackage{lmodern}
\usepackage{microtype}
\usepackage{amsmath, amssymb, amsthm}
\usepackage{geometry}
\geometry{a4paper, margin=1in}
\usepackage{hyperref}
\hypersetup{colorlinks=true, linkcolor=blue!50!black, citecolor=blue!50!black, urlcolor=blue!50!black}
\usepackage[backend=biber, style=numeric-comp, sorting=nyt, natbib=true, defernumbers=true]{biblatex}
\BiblatexSplitbibDefernumbersWarningOff
\addbibresource{intelligence_compressors.bib}
\usepackage{xcolor}

\title{Intelligence Is the Discovery of Compressors}
\author{Vasili Gavrilov\\
\small Independent researcher\\
\small ORCID: \href{https://orcid.org/0009-0007-9371-5994}{0009-0007-9371-5994}}
\date{\today}

\begin{document}
\maketitle

\begin{abstract}
Intelligence is the discovery of compressors whose inductive bias matches the structure of the data they encounter. Not every compression is an act of intelligence, but every compressor was found by one, and the discriminator is the distinction between the act and the artifact. Inventing the Fourier transform, the Huffman code, or least-squares regression was an act of intelligence; running the resulting algorithm over a buffer replays a search that was completed once and has been amortized across every invocation since. Between the two poles lies the case that recurs daily: judging which known compressor fits the data at hand, which is discovery at reduced amplitude and intelligent in the same sense, in smaller measure. Intelligence in compression is therefore a gradient of amortization, not a binary. The bias-variance axis names the two cheap failure modes on either side of the useful band: overfitting (memorization of specific examples without generative structure) and underfitting (over-smoothing that loses the structure); what lies between them is named by the word both are built from, \emph{fitting}, and it is an operation rather than a location, the iterative bracketing of an unknown structure from the too-specific and the too-general side until the description is as short as the structure allows and no shorter. Where the bracket closes is set by the observer's threshold for ``fits well enough'': beauty in the eye of the beholder, intelligence in the eye of the beholder. We connect four formal traditions of compression (Solomonoff induction and Kolmogorov complexity, the Information Bottleneck, the Renormalization Group, and weight-decay-as-complexity-prior) and locate compressor discovery at three substrate levels: (a) evolutionary search across generations (genetic and cultural inheritance); (b) researcher choice and orchestration of tooling (methods, architectures, pipelines); (c) self-organizing dynamics (multi-agent self-play and self-organizing networks); each paying the search cost in a distinct currency. Empowerment maximization supplies the motivational layer linking compression to evolutionary fitness; as a structural consequence, sufficiently agentic systems pursue empowerment as an instrumental sub-goal regardless of substrate. Contributions: (i) three refinements that separate an act of intelligence from the execution of its product, namely the act/artifact distinction with its intermediate case, the fitting band read as two-sided iterative approximation, and the observer-relative threshold; (ii) the three-substrate-level taxonomy of compressor discovery; (iii) empowerment as the motivational complement to compression-as-mechanism. Grounded in three decades of hands-on work across tree search, expert systems, genetic algorithms, and feed-forward neural networks.
\end{abstract}

\section{Introduction}
% PURPOSE: State the thesis (intelligence = discovery of compressors); preview
% the structure; scope the contribution honestly as synthesis rather than
% original formal result.

Different research traditions in artificial intelligence, machine learning, statistical physics, and theoretical neuroscience have, often in isolation, arrived at structurally similar claims about what intelligence is and where it comes from. Information theory \citep{shannon1948} says intelligence is compression: a system that predicts well has, of necessity, found regularities in its inputs and represented them compactly. Algorithmic information theory says the same in stricter language: the shortest program generating the data is the data's intrinsic structure, and any system that discovers it has discovered the data's content \citep{solomonoff1964, kolmogorov1965}. Statistical physics says it in the language of renormalization: useful descriptions arise by integrating out detail that does not matter at the scale of interest. Theoretical neuroscience, via the Free Energy Principle and predictive processing, says that organisms maintain themselves by becoming compressors of their environments. Modern deep learning has demonstrated that weight regularization corresponds to a complexity prior \citep{hinton_vancamp_1993}, and the architectural inductive biases of convolutional \citep{lecun1989} and transformer networks are themselves compressed encodings of the structural assumptions a useful learner needs.

These traditions converge on a family of claims but do not name the convergence, and the naive convergent claim is too permissive: a great many algorithms compress while doing none of the work intelligence does, because the work was done once, by whoever found them, and is thereafter only replayed. We argue that intelligence is the discovery of compressors whose inductive bias matches the data, judged against an observer-relative threshold, with overfitting and underfitting as the cheap failure modes on either side of the band we name \emph{fitting} (developed in Section~\ref{sec:discriminator}). The thesis has direct precedent in \citet{hutter2005} (AIXI, the Solomonoff prior); the act/artifact discriminator with its intermediate case, the fitting and observer-relativity refinements, the three-substrate-level taxonomy, and the pairing with empowerment-as-motivation are this paper's contributions. The apparent diversity of intelligent systems reflects different ways of paying the search cost for such compressors, and the substrate is incidental to the discovery itself \citep{gavrilov_emergence_2026}.

Three substrate-level mechanisms account for compressor discovery observed in practice. (a) Evolutionary search across generations (both genetic and cultural inheritance: language, pedagogy, traditions) pays the cost in fitness differentials over evolutionary timescales. (b) Researcher choice and orchestration of tooling (methods, architectures, pipelines) compresses by encoding domain priors into the structure of the system or the choice and composition of inferential methods, paying the cost in researcher effort. (c) Self-organizing dynamics (covering competitive multi-agent self-play and self-organizing networks like Kohonen, Hopfield, and Boltzmann machines) discover emergent compressors through component interaction under local rules, paying the cost in compute. The three cover the canonical instances of compressor discovery.

Compression accounts for the mechanism of intelligence but not its motivation. The motivational layer is supplied by the empowerment and causal-entropic-force tradition: organisms persist by maintaining access to many possible future configurations, and compressors are valuable because they enable that future option-space \citep{wissner_gross_2013}. Compression and empowerment are two faces of the same underlying selection pressure. Better compressors enable better predictions; better predictions enable better actions; better actions enlarge the agent's future state space; and an enlarged future state space is the fundamental currency of evolutionary survival. The framework predicts that any system intelligent enough to act over long planning horizons will pursue empowerment as an instrumental sub-goal, regardless of substrate or terminal goals \citep{bostrom2014}.

The contribution is the unified framing plus the act/artifact, fitting, and observer-relativity refinements, not the constituent ideas in isolation. Several questions remain deliberately bracketed: the phenomenology of subjective experience and the precise relation between compression and out-of-distribution generalization. The framework should be judged on whether it makes the convergence between traditions easier to see, not on whether it resolves their open problems.

The author is an independent researcher with roughly three decades of industry implementation experience in AI-adjacent and decision-system engineering. Production-deployed work includes: a linear-regression classifier for automatic medical-case assignment; a rule-engine subsystem for health-care finance and statistics; a genetic-algorithm labels-placement system for commercial charting and plotting; predictive analytics in health data science; tooling for statistical (noisy) genome matching; search-engine programming; and projects in entity identification spanning compression, encryption, and search. The underlying engineering substrate has been daily work with relational databases, i.e. relational algebra over sets, where every well-formed query is itself a compression of source relations (the data) into the structure of the data of interest via selection, projection, join, and aggregation. Educational background spans undergraduate physics (1991 programming work on scientific and medical nuclear-magnetic-resonance apparatus, involving Fourier-domain spectroscopic compression) and computer science (1997 seminar thesis on backpropagation feed-forward neural networks). Earlier informal articulations and partial implementations include AIS (a personal-use indexing implementation from 2013) and Context Renormalization \citep{gavrilov_context_renorm_2026}, a 2025 protocol released on GitHub. Time-stamped public artifacts (1997 thesis, 2001 SMBPANN design notes, 2013 AIS, 2025 Context Renormalization, 2026 Substack essays cited throughout) serve as priority claims, not substitutes for peer review.

The remainder of the paper is organized as follows. Section~\ref{sec:observations} establishes the central empirical observation through three personal implementations (1995--2001) that appeared intelligent to external observers without source access, but not to the author who wrote every line. Section~\ref{sec:discriminator} refines the compression thesis by separating the act of finding a compressor from the artifact that results, naming the band between the two failure modes \emph{fitting} and reading it as a two-sided iteration rather than a place on an axis, and making the threshold for ``fits well enough'' observer-relative. Section~\ref{sec:heuristics} reframes heuristics in neural networks as inductive bias. Section~\ref{sec:discovery} follows the trajectory from hand-coded to discovered heuristics via the AlphaGo and AlphaZero lineage. Section~\ref{sec:formal} develops the formal foundation: the three entropies and the bridges connecting them. Section~\ref{sec:substrates} lays out the three substrate-level mechanisms of compressor discovery. Section~\ref{sec:substrate_independence} distinguishes ontological substrate-independence from operational substrate-dependence. Section~\ref{sec:beyond_skull} extends the compression frame beyond the biological substrate \citep{clark_chalmers_1998}. Section~\ref{sec:empowerment} returns to the motivational question via empowerment and causal entropic forces. Section~\ref{sec:implications} traces implications for current artificial-intelligence research. Section~\ref{sec:conclusion} concludes.

\section{Three Empirical Observations: Minimax, Prolog, Genetic Algorithms}\label{sec:observations}
% PURPOSE: Establish via three hands-on implementations the perception of intelligence is observer-relative and paradigm-independent. Three substrates (tree search, symbolic inference, population evolution); one perceptual phenomenon. Formal mechanism: sensitive dependence in non-linear dynamics (Arnold ODE) and discrete-rule equivalent (Wolfram). Mirrored by Gilovich's catalog of human self-mis-attribution. The phenomenon is one perceptual symmetry, not two separate puzzles.

The central empirical observation: the perception of intelligence is observer-relative and paradigm-independent. Between 1995 and 2001 the author implemented three algorithmic systems on three different substrates. In each case, the system appeared intelligent to external observers without access to the source. To the author, who had written every line, the appearance was attenuated but not absent: playing his own creation in real time, especially when its next move could not be predicted, the impression of interacting with an intelligent program persisted. The effect survives partial transparency: it arises from the gap between system behavior and observer predictive capacity, not from ignorance of the implementation alone. The observation generalizes across the three paradigms, which suggests the effect is structural rather than tied to any particular algorithmic family.

The first system was a minimax search engine for checkers and chess in C, 1995 undergraduate coursework. The implementation followed classical minimax with alpha-beta pruning and a hand-coded position-evaluation function. At depths of four to six plies, the program defeated non-expert opponents more often than not. Observers without source access called it ``thinking''; to the author, who had written the programs himself, it was a recursive tree-search routine with a value-counting function at the leaves. To a player who could not see the implementation, it looked like a thinking entity.

The second system was an expert system in LPA Prolog, 1996 undergraduate coursework, drawing on the AI-Prolog idiom established by \citet{bratko1990}. The program performed rule-based inference over a knowledge base of declarative facts and rules, producing diagnostic conclusions from observed evidence in a narrow domain. Prolog's symbolic-inference paradigm is structurally quite different from minimax tree search: no game tree, no leaf-evaluation function, no alpha-beta pruning. Yet the same observer-relative effect appeared. From outside, the diagnostic conclusions had the quality of considered judgment. From inside, the program was depth-first symbolic backtracking over a fixed set of rules. Two paradigms; one perceptual effect.

The third system was a genetic-algorithm-based optimizer (2001) for a commercial spatial-optimization problem: labels-placement in a charting and plotting framework. Unlike the prior two systems, the GA had no explicit search tree and no fixed inference procedure; it maintained a population of candidate solutions, applied mutation and crossover operators, and selected for fitness across generations. The system found near-optimal placements where exhaustive search was infeasible. Observers described the convergence in terms suggestive of agency (``the algorithm figures out where to put the labels''); to the author, who had written the mutation operators and the fitness function, it was a population-based stochastic optimizer converging on a fitness landscape.

Three paradigms (tree search, symbolic inference, population evolution); three implementations; one perceptual phenomenon. The pattern is paradigm-independent. The effect depends not on algorithmic family but on the observer's ability to predict the system's next move in real time, which each of these systems outruns.

The formal mechanism is well-understood in the dynamical-systems literature: deterministic systems generate trajectories that are practically unpredictable in finite time when they are non-linear and sensitive to initial conditions \citep{arnold_ode}. A double pendulum is deterministic but practically chaotic; the Lorenz attractor is fully specified by three ODEs yet produces trajectories that diverge exponentially under any finite-precision approximation. Minimax over a sufficiently large game tree, Prolog's symbolic search, and GA dynamics are all deterministic systems whose trajectories are, from the observer's perspective, identically unpredictable. \citet{wolfram2002} makes the discrete-rule version: simple deterministic rules over cellular automata generate macroscopic patterns indistinguishable from those produced by intelligent agency until the rule itself is inspected. The ``looks intelligent'' effect is what we should expect from any deterministic system whose trajectory is computationally rich enough to outrun the observer's predictive capacity.

The mirror observation is that humans frequently attribute intelligence to themselves where the underlying mechanism is closer to pattern completion than deliberative reasoning. \citet{gilovich1991} catalogs the systematic biases by which human inferential behavior departs from normative judgment: confirmation bias, illusory correlation, base-rate neglect, regression-to-the-mean blind spots, attribution errors. As algorithms look more intelligent than they are, humans look more intelligent to themselves than they are. The two are the same phenomenon: the gap between behavior and the observer's predictive capacity, whether the observer is external (looking at an algorithm) or internal (looking at one's own reasoning). The boundary between ``intelligent'' and ``algorithmic'' is not where ordinary observation places it.

This sets up the question for the rest of the paper. If the perception of intelligence is observer-relative and paradigm-independent, what is the actual structural property that distinguishes intelligent systems? Section~\ref{sec:discriminator} develops the answer.

\section{The Act and the Artifact: Invention, Application, and the Fitting Band}\label{sec:discriminator}
% PURPOSE: Refine the thesis. The discriminator is act vs artifact, not compressor-type. Invention is intelligent; execution is amortized replay; matching a known compressor to data is the intermediate case. Bias-variance failure modes and the fitting band as two-sided iteration. Cost asymmetry on serial hardware. Observer-relativity of the threshold. "Intelligence in the eye of the beholder."

The compression thesis as stated so far is over-permissive, but the usual way of narrowing it is wrong. It is tempting to say that linear regression, the Fourier transform, and MP3 audio compression all compress without being intelligent, and to file them as counterexamples. They are not counterexamples. They are the thesis with its history deleted. Every compressor in that list was found by an act of intelligence: Fourier's decomposition of an arbitrary signal into sinusoids, and later the reduction of its cost from $O(n^2)$ to $O(n\log n)$ that made it usable; Huffman's optimal prefix code; the dictionary construction of Ziv and Lempel that sits underneath every general-purpose archiver in daily use; least squares. Each was a discovery of precisely the kind this paper says intelligence consists of, each was hard, and each arrived once. What is not intelligent is the \emph{execution}. Running an FFT over a buffer, or a dictionary coder over a file, replays a search that was completed in a mind and has been amortized ever since across every subsequent invocation. The confusion is between the act and the artifact: an algorithm is the fossil of an act of intelligence, not an instance of one, and that the fossil is cheap to copy is exactly what made the original act worth so much.

Between invention and execution lies the case that occurs daily, and it is the one worth being careful about. Applying a compressor requires deciding that it applies: that this signal is sparse in the frequency domain and not the time domain, that a wavelet basis suits a feature localized in both, that this relation is linear in these covariates and not those, that this corpus carries the repetition statistics a dictionary coder exploits. None of that is contained in the algorithm, and the data does not announce it. It is a match between an inductive bias held in advance and a structure encountered, which is the same operation as discovery, conducted over a small catalogue of known compressors rather than an open space of possible ones. Using a compressor as an entropy reducer is therefore intelligent in the same sense as inventing one, in smaller measure, and the measure is set by how much of the search a prior act has already paid for. Intelligence in compression is a gradient of amortization rather than a binary, and substrate (b) of Section~\ref{sec:substrates}, researcher choice and orchestration of tooling, is this activity given its own name; that it is a genuine mode of compressor discovery, and not a lesser administrative step beneath the other two, follows directly.

Section~\ref{sec:observations} reads differently in this light. The author found his own minimax, Prolog, and genetic-algorithm systems less intelligent than external observers did, and unpredictability is only half the explanation. The other half is that he had performed the intelligent part himself, in advance: the position-evaluation function, the rule base, the fitness function and the mutation operators. Seen from inside, each running program was the replay of his own completed search. Seen from outside, with the search invisible, replay and act were indistinguishable. Observer-relativity and the act/artifact distinction are two descriptions of one phenomenon, and the rest of this section develops both: a structural filter (the bias-variance axis from classical statistics) and an evaluation criterion (an observer-relative threshold).

The bias-variance axis names the two cheap failure modes that lie on either side of the intelligent band. \emph{Overfitting} memorizes specific examples without recovering the generative structure of the data: the compressor stores instances rather than the pattern that produced them. \emph{Underfitting} smooths the signal so heavily that the relevant structure is lost: the compressor produces a low-order summary that misses what mattered. Both failure modes produce compressed representations; neither produces intelligence. The author has elsewhere catalogued the consequences in current artificial-intelligence use without a strong prior given \citep{gavrilov_machines_generalize_2026}: large language model (LLM) tools default to overfitting (memorizing training-set patterns as if they were universal) and to underfitting (importing generic enterprise-architecture complexity where the structure of the specific problem requires the opposite). The intelligent band is the expensive middle: a compressor whose inductive bias matches the structure of the data, neither memorizing examples nor smoothing them away.

Naming that middle band is worth doing carefully, because the plainest available word is also the most accurate. If the failure on one side is overfitting and the failure on the other is underfitting, then what lies between them is \emph{fitting}, in the ordinary sense of the word: fitting a curve to points, fitting an objective function to a problem, fitting a key to a lock. The word is normally reserved for the mechanical part of statistics, the part a solver does after the modelling is finished; the claim here is that it names the whole operation. And it names an operation, not a location. No procedure jumps to the correct complexity in one step, because the correct complexity is a property of a structure not yet known; what procedures do instead is approach it from both sides and tighten. A network is initialized too crude to represent the data and trained toward the overfit regime until held-out error turns, so the stopping point is bracketed rather than computed. A model class is swept from too simple to too complex and the elbow is read off the curve. A theory is proposed too general, corrected against anomalies until it is nearly too specific to the cases that produced it, then generalized again. Regularization is the formal name for pulling back from the overfit side and added capacity for pushing up from the underfit side, and the two are run against each other until the residual stops improving. Intelligence, on this reading, is successive approximation of an unknown structure by bracketing it between the too-specific and the too-general, and the compressed description is what the bracket converges to. This is also why the middle is expensive while both failures are cheap: each failure is a single move, and fitting is an iteration. Where the bracket is allowed to close, that is, how much fit is enough, is not fixed by the data either, which is the subject of the paragraph after next.

The cost asymmetry between the two failure modes and the intelligent band is severe and structurally important. Overfitting is computationally cheap: $O(1)$ lookup against a stored example. Underfitting is also cheap: omit the structure and predict a low-order summary. The intelligent middle requires real computation, and on serial von Neumann hardware the gap to the cheap shortcuts is brutal. This is why most ``intelligence'' failures in machine systems default to memorization rather than to genuine compression, and why the discovery of a \emph{fitting} compressor is the hard problem the three substrate-level mechanisms (Section~\ref{sec:substrates}) all confront. Massively parallel substrates, biological and otherwise, narrow the gap but do not close it; the structural asymmetry between lookup and computation persists. NP-hard combinatorial problems that take exponential time on serial hardware can be approximated in polynomial time, or faster, on parallel biological substrates that run physical dynamics whose fixed points encode the solutions (Section~\ref{sec:substrate_independence}). The author's 1997 thesis \citep{gavrilov_1997} framed this observation through Flynn's now-largely-forgotten taxonomy of parallel architectures: SISD (Single Instruction, Single Data: classical von Neumann), SIMD (Single Instruction, Multiple Data: vector and GPU operations), MISD (Multiple Instruction, Single Data: rare in practice), and MIMD (Multiple Instruction, Multiple Data: multicore CPUs, distributed clusters, biological substrates). A problem tractable on one Flynn class may be intractable on another.

The evaluator's threshold for ``fits well enough'' is observer-relative. The same Mozart sonata is profound to a trained ear and boring to an untrained one. The same minimax engine looks like thinking to a non-implementer and like a recursive function with $O(1)$ leaf-cost to the implementer. The three systems of Section~\ref{sec:observations} all sit in the intelligent band for some observers and outside it for others. Beauty is in the eye of the beholder; intelligence is too. The same observer-dependent structure appears in quantum mechanics: before observation, a quantum system can be in many possible states at once (a superposition); the act of measurement by an observer is what selects one definite state. Reality, like intelligence, is partly constituted by the act of observing it. The threshold is set by the observer's own prior compressors: above some level of fit between the system's output and the observer's evaluation criterion, the system passes the threshold; below it, the system is judged trivial. This is the framework's central feature, not a defect: an absolute scale could not explain why the same artifact is seen as intelligent or trivial depending on who is judging.

A separate disambiguation is worth flagging: intelligence is not consciousness, and the two are frequently conflated. The operational mechanism most often associated with consciousness is a periodic monitoring loop: the agent samples its environment, checks its state against a set of rules, and updates its behavior. Game programmers know this loop as the main game loop; robotics knows it as the sensorimotor cycle; \citet{hofstadter1979} mentioned a hypothesized brain region whose activity correlates with conscious states (whether subsequent neuroimaging has confirmed this specifically is an open empirical question the present paper does not adjudicate). Implementing such a loop is straightforward in both non-intelligent automata and intelligent systems. Consciousness-as-loop is orthogonal to intelligence-as-compression: the loop determines moment-to-moment awareness, while the compressor inside the loop determines whether the resulting behavior is intelligent. An intelligent system with a monitoring loop produces unpredictable-but-interesting behavior that meets our expectations of intelligence; a non-intelligent automaton with the same loop produces predictable scripted behavior. The present framework addresses only the compressor side and brackets the consciousness-loop as separate machinery that may or may not co-occur with intelligence.

The next section examines how heuristics encode the prior structure that pushes a learner toward proper fit and away from both failure modes.

\section{Heuristics as Inductive Bias}\label{sec:heuristics}
% PURPOSE: Reframe heuristics as compressed encodings of structural assumptions about the data. The 1997 thesis claim and its 25-year empirical confirmation. Bionics as the framing the author used at the time. LeCun's CNN as canonical bionics. Brain-architecture extrapolation (connectome, brain imaging). Connection to substrate-level (c). Operational form via the generalize-vs-specialize tension.

A core observation in the author's 1997 undergraduate thesis on backpropagation feed-forward neural networks was that the network's architecture and preprocessing pipeline carry as much of the learning work as the gradient-descent algorithm that adjusts the weights \citep{gavrilov_1997}. The line of thought began with a hint from the supervisor, Professor Ehud Gudes (to whom the author records his gratitude here for the suggestion), to investigate the role of heuristics in neural-network design; this opened the connection between bionics (biological design principles, developed below) and architectural choice as a necessity rather than an optimization. The hardware constraint was concrete: a typical mid-1997 desktop was the Pentium MMX (166--233 MHz, 16--32 MB RAM), and fully-connected networks of any meaningful size were computationally infeasible to train. Architectural compression of the hypothesis space was a precondition for any working system, not a refinement of it. At the time, the dominant framing treated backpropagation \citep{rumelhart1986} as the algorithmic core and network topology as a relatively minor engineering decision. The 1997 thesis disagreed: without problem-specific architectural commitments and preprocessing that compressed the input into a form the network could exploit, even a correctly-trained network would fail to generalize. The intervening twenty-five years have been an extended empirical confirmation. Convolutional neural networks (CNNs) \citep{lecun1989}, residual networks, attention-based transformers, and the foundation-model era have all shown that the architectural prior is at least as important as the optimizer. The optimizer is generic; the architecture is where the structure of the problem domain lives.

In compression terms: a neural network with a fully-connected universal-approximator topology can, in principle, learn any continuous function on a compact domain. In practice, this universal-approximation guarantee is operationally useless for finite training data: the hypothesis space is too large, and most functions consistent with a finite training set generalize badly to unseen data. The role of architectural inductive bias is to compress the hypothesis space, in advance of any training, down to the subspace where good solutions are likely to live. A convolutional network has dramatically fewer parameters than a fully-connected network of comparable expressive power because the convolutional structure encodes the prior that nearby pixels are statistically related and that the same feature detector is useful at different image locations. That prior is correct for natural images, which is why CNNs work on natural images and would not work on scrambled-pixel inputs without it. The architectural commitment is not an optimization trick; it is a compressed encoding of structural assumptions about the data \citep{russell_norvig_1995, bigus1998}.

The author's personal mental model for this position connected to bionics, a term he had encountered in a few popular-science articles read in high school in the 1980s; the 1997 thesis itself did not use the word, but the bionics connection resonated for the author and shaped the architectural-priors argument from outside the formal write-up. (Bionics covers what is now called biomimetics and bio-inspired engineering in current Western terminology, i.e. the practice of copying biological design principles for artificial systems.) LeCun's convolutional network is a clean instance of bionics applied to vision: the convolutional structure copies the local feature-detection architecture of the retina and early visual cortex, applying local receptive-field operations tiled across the image rather than fully-connected processing across the entire visual field at every pixel. The 1997 thesis pointed in this direction as its central prediction: continued copying of brain architecture into neural-network architecture would address the three coupled problems the field then faced: overfitting (too many parameters for too little data), underfitting (architectures too generic to capture the structure), and intractability (compute too limited to train fully-connected networks at any meaningful scale). The intervening decades have validated the prediction: biologically-inspired deep learning, neuromorphic computing, and large-scale brain-inspired architectures all inherit the same bionics intuition.

Bionics also makes the connection to the three-substrate-level taxonomy of Section~\ref{sec:substrates} legible. Researcher choice and orchestration of tooling is one of the three canonical modes of compressor discovery: human researchers borrow structural priors from biology because evolution has already paid the search cost. Substrate-level (b) is therefore not in tension with substrate-level (a) but borrows from it, and bionics is the engineering practice connecting the two.

Gradient descent itself predates the neural-network specific development: the classical numerical-methods literature on iterative descent against an objective function is the foundation on which backpropagation rests \citep{mudrov1991}. The novelty of backpropagation was not gradient descent itself but the chain-rule derivation that makes the gradient computable through deeply nested compositions of differentiable functions \citep{rumelhart1986}. The choice of gradient descent as the optimization method (over Kohonen self-organization, EM, MCMC, or any other iterative procedure) is itself a substrate (b) commitment; the within-method execution that follows pays the cost in compute and training data rather than researcher effort.

Inductive bias is therefore the prior that pushes a learner toward the intelligent middle band of Section~\ref{sec:discriminator} and away from both overfitting (memorize what should generalize) and underfitting (smooth away what should specialize). A compressor that finds short representations without the right prior about which short representations matter for the task finds the wrong ones. The same generalize-versus-specialize tension appears across domains: in software engineering with LLM tools, where the model defaults to enterprise-architecture complexity where the project is intentionally simplifying \citep{gavrilov_machines_generalize_2026}; in medicine, where personalized treatment outperforms generic protocols but cannot be applied at full scale due to cost; and in hardware, where specialized tooling outperforms general-purpose multitools at specific tasks but cannot replace them everywhere. The general-purpose-versus-specific-purpose trade-off is the operational consequence of asking a compressor to find regularities without telling it which regularities matter.

If those commitments can be hand-coded by researchers, can they also be discovered by competitive search rather than written by humans? Section~\ref{sec:discovery} examines the AlphaGo and AlphaZero lineage as the cleanest demonstration.

\section{From Hand-Coded to Discovered Heuristics: From AlphaGo to AlphaZero}\label{sec:discovery}
% PURPOSE: Heuristics can be FOUND by competitive search rather than written by humans. AlphaGo (2016) -> AlphaGo Zero (2017) -> AlphaZero (2017) lineage. Nuance: game-specific heuristics were discovered, but architectural bias (rules, network architecture, MCTS, reward) was still hand-engineered. The substrate-level taxonomy of Section 6 picks this distinction up.

Heuristics need not be supplied by humans. The AlphaGo--AlphaZero line at DeepMind (2016--2017) demonstrated this most cleanly: AlphaGo (2016, with human-game bootstrapping), AlphaGo Zero (2017, from random initialization via self-play), and AlphaZero (2017, generalized to chess and shogi from random initialization) successively eliminated human-supplied heuristics while still defeating the strongest classical engines. Heuristics that chess engines had accumulated over decades were rediscovered from self-play alone.

The author's 2001 genetic-algorithm labels-placement system (Section~\ref{sec:observations}) exhibited the same principle earlier and at smaller scale: no human-supplied heuristics for where each label should go; the competing generations discovered near-optimal placements over a few generations, while only the mutation operators, crossover scheme, and fitness function (substrate (b) work) were hand-engineered.

What AlphaZero replaced and what it did not is the load-bearing distinction. Domain-specific heuristics were discovered: piece-square tables, opening preferences, position-evaluation functions. Non-domain-specific inductive bias remained hand-engineered: game rules, network architecture, search structure (Monte Carlo Tree Search, or MCTS), reward signal, training procedure. Discovery happened at one substrate level; hand-engineering remained at the level above it. Section~\ref{sec:substrates} taxonomizes the two as distinct modes of compressor discovery, and the rest of the paper turns on this distinction. Section~\ref{sec:formal} first develops the formal apparatus underneath both.

\section{Formal Foundation: Three Entropies and Their Bridges}\label{sec:formal}
% PURPOSE: Lay out the formal apparatus. Three entropies (thermodynamic, informational, algorithmic) connected through Jaynes-style max-entropy. Solomonoff. Information Bottleneck (narratively). Renormalization Group. Weight-decay as complexity prior (Hinton-Van Camp). Gradient descent operationalizes the principle.

Compression has formal expression in three traditions whose central quantities are technically distinct but operationally aligned: statistical mechanics, information theory, and algorithmic information theory. Each defines an entropy; the three are connected by bridges that are mathematically tight in restricted settings and conceptually unified in general.

\textbf{Boltzmann's thermodynamic entropy}:
\begin{equation}
S = k_B \ln W,
\end{equation}
where $W$ is the number of microstates compatible with a given macrostate.

\textbf{Shannon's informational entropy} \citep{shannon1948}:
\begin{equation}
H = -\sum_i p_i \log p_i,
\end{equation}
the expected number of bits required to specify the outcome of a random variable distributed as $p$.

\textbf{Kolmogorov complexity} \citep{kolmogorov1965}:
\begin{equation}
K(x) = \min_p \{|p| : U(p) = x\},
\end{equation}
the length of the shortest program for a universal Turing machine $U$ that outputs the string $x$.

These three quantities count, respectively, microscopic configurations, probabilistic outcomes, and symbolic descriptions. They are not the same quantity, but each measures the irreducible information content of a state, a distribution, or a string.

The bridge from Boltzmann to Shannon is Jaynes's maximum-entropy formulation of statistical mechanics. Imagine measuring a gas's macroscopic properties (temperature, pressure, volume) without seeing the microscopic arrangement of atoms. Among all probability distributions over the microscopic states consistent with the measurements, the one with the maximum uncertainty (Shannon entropy) turns out to give exactly Boltzmann's thermodynamic entropy. Statistical mechanics, in this reading, is information theory applied to physical systems whose only access is via macroscopic measurements. The bridge from Shannon to Kolmogorov: for long enough sequences of symbols generated by a computer program, the Shannon entropy per symbol (the average uncertainty) and the Kolmogorov complexity per symbol (the average shortest-program length) converge to the same value as the sequence grows. The three entropies are different formalizations of how much information is needed to specify the state.

Solomonoff induction connects Kolmogorov complexity to learning. The universal prior \citep{solomonoff1964} assigns to each computable hypothesis $h$ a probability inversely related to its description length:
\begin{equation}
P(h) \propto 2^{-K(h)}.
\end{equation}
A Bayesian agent using this prior is asymptotically Bayes-optimal in a strong sense: any reasonable prior converges to it. The principle is not directly computable, since Kolmogorov complexity is uncomputable, but it formalizes the preference for shorter explanations of the data. Practical learning algorithms are computable approximations to it.

The Information Bottleneck (Tishby and collaborators) and the Renormalization Group are two specializations of compress-by-discarding-what-does-not-matter, differing on what defines mattering. The Information Bottleneck targets task relevance: given input $X$ and target $Y$, it seeks a representation $T$ that throws away as much detail from the input as possible while keeping everything that matters for predicting the target. Formally, this means minimizing the mutual information $I(X;T)$ (how much $T$ remembers about $X$) while preserving $I(T;Y)$ (how much $T$ tells us about $Y$). Each layer of a trained deep neural network has been argued to approximate one step of this trade-off, though the exact relationship is debated. The Renormalization Group targets scale. RG transformations zoom out from a system's small-scale fluctuations, averaging them away to produce a simpler effective description for large-scale behavior. The striking result is universality: many systems with entirely different microscopic physics converge to the same effective theory at large scales. Water near its boiling point and magnets near their critical temperature, for example, show identical fluctuation behavior at large scales despite being completely different at the atomic level.

Signal processing offers a different operational form of compression: the Fourier representation. A signal that is dense in the time domain is typically sparse in the frequency domain; the Fourier transform performs the basis change that exposes the sparsity. The same principle underlies spectroscopy, nuclear-magnetic-resonance tomography, JPEG image compression (via the related Discrete Cosine Transform), and signal-analysis applications generally; wavelets generalize the construction to representations matched to localized features. In each case, compression is the recognition of the right basis for the data: the basis in which the signal is short. The recognition is the intelligent part; the transform, once written down, is machinery. Constructing these bases was an act of the first kind, and every subsequent application is an act of the second, less the residual judgement of which basis to reach for, which is the intermediate case of Section~\ref{sec:discriminator}.

Modern neural-network training enacts a related principle: regularizing the magnitude of the network's weights selects, among the parameterizations that fit the training data, those with shorter description length under a fixed-precision encoding \citep{hinton_vancamp_1993}. Weight decay, originally introduced as an optimization trick to reduce overfitting, is in this reading a soft enforcement of Solomonoff's preference for shorter explanations. The underlying minimum-description-length principle was already explicit in the 1993 formulation.

Gradient descent is the procedure that finds these short representations in practice: iterative numerical optimization \citep{mudrov1991} descending the loss surface of the parameterized hypothesis class. Combined with weight regularization, it operationalizes the abstract compression principle on physical hardware. The formal apparatus reviewed above specifies what gradient descent is, in effect, searching for. The next section asks where in the engineering stack that search can occur.

\section{Compressor Discovery at Three Substrate Levels}\label{sec:substrates}
% PURPOSE: Central organizing claim. Three modes of compressor discovery, paying search cost on different substrates.

Compressor discovery occurs in at least three canonical modes. Each pays the search cost on a different substrate.

\textbf{(a) Evolutionary search across generations.} The substrate is the genome and the cultural-transmission machinery; the cost is paid in fitness differentials over evolutionary timescales; the result is inherited at birth and at acculturation. Organisms across all scales of life (bacteria, plants, animals), their constituent systems (eyes, brains, immune systems), the architecture of biological cognition, and the cumulative inventions transmitted across human generations are all products of this search. Cultural inheritance runs in parallel as a roughly equally significant extra-genetic transmission channel: language, pedagogy, and traditions (religious, professional, scientific) all pass compressors from one generation to the next without each individual re-discovering them. Nature and nurture are not cleanly separable, and both shape the inherited prior. Substrate (a) is slow, expensive in lives, and globally parallel. The discovered compressors are inscribed in the genetic, developmental, and cultural machinery of organisms: an organism does not learn them, it inherits them. Bionics (Section~\ref{sec:heuristics}) is the engineering practice of borrowing the results.

\textbf{(b) Researcher choice and orchestration of tooling.} The substrate is human research and design effort; the cost is paid in researcher time; the result is a choice of methodology, architecture, and tools (and increasingly, an orchestration of multiple such choices into a pipeline) that compresses the hypothesis space in advance of any data being processed. Substrate (b) is where the researcher selects which inductive bias to bring to a problem. Outputs of (b) include: architectural priors (convolutional networks \citep{lecun1989}, attention mechanisms, residual connections); classical statistical methods (least-squares / MSE regression, K-means clustering, principal component analysis); broader methodological commitments (Bayesian vs frequentist, Markov Chain Monte Carlo (MCMC) vs variational inference, symbolic search vs gradient-based optimization); and orchestration of multiple tools (an LLM is itself a substrate (b) output, but the agentic pipeline that composes an LLM with retrieval, other models, and tool calls is substrate (b) work at a meta level). Substrate (b) borrows from (a): the researcher takes structural ideas from biology (where evolution has already paid the search cost) and methodological ideas from the accumulated statistical and computational canon (where culture has). Bionics is the explicit form of architectural borrowing; the established statistical canon is the explicit form of methodological borrowing. The within-method execution that runs after a (b) choice (gradient-descent iteration, EM convergence, K-means iteration, MCMC sampling) is the consequence of the (b) choice, not a separate substrate; for some tasks the execution discovers a lot (deep neural-network weight training), for others it is near-mechanical (K-means converging to local fixed points), and for yet others it is trivial (closed-form / analytical solutions like the least-squares normal equation).

\textbf{(c) Self-organizing dynamics: replicating substrate (a) in compute time.} The substrate is a system whose components interact under local rules, producing emergent global structure without explicit specification by the researcher. The cost is paid in compute (rather than the evolutionary timescales of substrate (a)). Substrate (c) is structurally equivalent to substrate (a) in mechanism (both produce emergent compressors through local competition and selection without global design) but replaces fitness differentials over generations with iteration over compute, making the same dynamics dramatically faster. Two sub-categories cover most canonical instances. \emph{Competitive multi-agent self-play}: AlphaZero (Section~\ref{sec:discovery}), generative adversarial networks, population-based training, and competitive reinforcement-learning regimes. In each, a population of agents competes under a shared fitness signal, with the curriculum self-organizing because opposition difficulty tracks the current population's strength. \emph{Self-organizing networks}: Kohonen self-organizing maps, Hopfield networks, Boltzmann machines, and similar systems where internal units compete or interact under local rules within a single network. In both sub-categories, substrate (c) discovers structure that no human explicitly wrote down, and is structurally distinct from substrate (b) because the discovery emerges from the dynamics rather than from researcher specification: the researcher sets up the game (or the topology and update rule) at substrate (b), but the discovery that follows comes from the dynamics themselves. The same self-organizing principle, applied at a meta level, can be used to search over substrate (b) choices: neural architecture search and population-based training are canonical applications, where substrate (c) is deployed to find the most-fit substrate (b) configuration for a given problem.

The three are not alternatives but layers of one stack. AlphaZero illustrates this directly: substrate (c) discovered the chess-specific heuristics described in Section~\ref{sec:discovery} while resting on substrate (b) commitments (architecture, search structure, reward signal) and ultimately on substrate (a) (the neural-systems intuitions behind value-and-policy networks, which biology and the accumulated research culture provided). Lowering the cost at one level shifts the bottleneck upward, where the next generation of research effort then concentrates. The engineering choice is which level to absorb the cost at, and increasingly, how to compose substrates within a single pipeline. Section~\ref{sec:substrate_independence} examines what stays constant across substrates and what does not.

\section{Substrate Independence: Principle and Limits}\label{sec:substrate_independence}
% PURPOSE: Distinguish ontological from operational substrate-independence. Both true simultaneously.

Substrate-independence holds in two distinct senses, one ontological and one operational. Conflating them produces most of the confusion in current debates.

The ontological claim is multi-realizability: the same computational and causal structure can be physically realized in many substrates. Intelligence is the property of a pattern, not of hardware. The author's prior essay \citep{gavrilov_emergence_2026} defends this directly: emergent phenomena arise from non-linear dynamics regardless of what is doing the dynamics. The substrate is the medium, not the message.

The empirical witness on the biological side is unambiguous. Insect superorganisms exhibit collective intelligence (foraging optimization, climate regulation, division of labor, nest architecture) on a substrate of distributed chemical signals among individually-simple agents, without anything resembling a central nervous system \citep{holldobler_wilson_2008}. The intelligence is at the colony level; the substrate is computationally and physically nothing like a neural network. \citet{bloom2000} generalizes the same observation across all scales of life, from bacteria to primates, on radically different substrates.

The formal mathematical foundation comes from non-linear dynamical systems. Structural stability and bifurcation analysis show that qualitatively new patterns (limit cycles, strange attractors, bifurcation cascades, pattern formation) arise from continuous parameter changes in ways not reducible to the generating equations \citep{arnold_catastrophe}. The same emergent structures appear across systems whose underlying equations differ substantially: the dynamical-systems mirror of multi-realizability. \citet{wolfram2002} makes the discrete-rule version: simple substrates produce structurally similar computational behaviors at the macroscopic level. Across these traditions the message is consistent. The macroscopic pattern carries the cognitive content; the microscopic substrate is incidental at the appropriate level of description.

The operational claim is distinct: different substrates have different efficiency profiles, different failure modes, and different costs of finding particular kinds of compressors. A biological brain finds face recognition trivial and 8-digit multiplication agonizing; a transistor-based computer has exactly the inverted profile. A rule-engine system finds policy violations easily but cannot do gradient-descent learning; a feed-forward neural network does gradient-descent learning easily but cannot easily encode hard logical constraints. The slime mold \emph{Physarum polycephalum} approximates solutions to hard network-optimization problems (specifically Steiner trees, the near-optimal way to connect a set of points) without any central nervous system. The mechanism is chemotaxis: cytoplasm preferentially flows toward food sources, and the network of cytoplasmic tubes reorganizes as material redistributes. These problems are NP-hard on serial von Neumann hardware, meaning no step-by-step algorithm solves them efficiently; the mold makes them tractable not by reasoning about them but by running physical dynamics whose stable end-states happen to be the solutions. Humans, modeling such problems symbolically, find them hard; the mold, not modeling them at all, finds them in proportion to its growth rate. The three substrate levels of Section~\ref{sec:substrates} are not interchangeable from an engineering standpoint. Each pays its discovery cost in a different currency, and the discovery itself happens in a different region of the compressor-space. Substrate (a) finds robust, slow, decentralized compressors over generations; substrate (b) finds well-understood compressors aligned with researcher intent and accumulated methodological canon; substrate (c) finds emergent compressors whose structure was not specified by any single tool: fast versions of the (a)-mechanism in compute time rather than evolutionary timescales, sometimes hard for humans to interpret.

Both claims hold simultaneously. Ontological substrate-independence permits intelligence across many substrates; operational substrate-dependence determines which substrate is appropriate for which task. Together they explain why we see intelligence across radically different physical implementations (ants, brains, networks) without those instances being equivalent in capability. Section~\ref{sec:beyond_skull} extends the compression frame across substrate boundaries, examining cognition that spans biological and engineered substrates simultaneously.

\section{Compression Beyond the Skull: Extended Cognition and Information-Exchange Protocols}\label{sec:beyond_skull}
% PURPOSE: Compression extends beyond biological substrate. Extended Mind framework. AIS + Context Renormalization. Industry parallels.

If compression is the mechanism of intelligence, it need not be confined to the biological organism. External tools that compress lifetime experience into searchable structure participate in the same cognitive operation. The philosophical framing is Extended Mind \citep{clark_chalmers_1998}: under the parity principle, a process that would count as cognitive when performed in the head still counts as cognitive when performed via external scaffolding meeting the conditions of reliable availability, automatic access, and trusted endorsement. The Otto-and-Inga thought experiment formalizes the move. Inga, with normal memory, recalls from biological memory that the museum is on 53rd Street and goes there. Otto, who has Alzheimer's, consults his constantly-maintained notebook to find that the museum is on 53rd Street and goes there. If we count Inga's biological recall as a case of memory, then by parity of reasoning Otto's notebook consultation is also a case of memory: the notebook is functioning as part of his mind.

The author's engineering work on this thesis spans two implemented systems separated by twelve years. AIS (Associative Indexing Service, 2013) was a personal-use indexing system designed to extend human associative memory: a personal Google over the user's own lifetime archive. It grouped sets and stored them as compressed entities retrievable by combinations of keys in any order, supporting the kind of associative recall the user's biological memory could no longer guarantee. The compression was pure (the same content stored under multiple key-combinations without duplication) and the retrieval was personalizable (the keys reflected the user's own associative structure rather than a generic ontology). It met the Clark-Chalmers conditions for the author who built and used it; it was abandoned when commercial mobile personal-knowledge-management tools subsumed the use case.

Context Renormalization \citep{gavrilov_context_renorm_2026} (2025, GitHub-released) is the successor system, targeting a different cognitive bottleneck: bounded large-language-model session continuity. It is a file-based protocol with an immutable header, an editable knowledge region, and a Compression Difficulty Score that signals when knowledge has outgrown the ledger and should migrate to durable artifacts. The protocol enforces a fixed character budget while preserving semantics under compression; when compression without loss becomes impossible, the system signals it explicitly rather than failing silently. The substrate is text on disk; the compressor is the human operator working with the protocol; the cognitive function performed is preservation of working knowledge across sessions, models, and compression.

The two systems together exemplify two layers of extended compression. AIS-class systems are personal compressors over lifetime data: the user's own memory extended through an external substrate. Context-Renormalization-class systems are protocols for compressed-concept exchange between minds and machines, skipping the natural-language re-encoding that ordinary conversation imposes. Sending a structured concept-graph between two agents is cognitively denser per unit bandwidth than sending the same content as prose, and the protocol layer is where that density lives.

Industry trajectory has validated the architecture. Anthropic's CLAUDE.md and OpenAI's memory layer function as project-level persistent compressed contexts; the SKILL.md format, converged on across major vendors in late 2025, standardizes capability-exchange. Vector databases and semantic embeddings implement compressed-concept exchange between AI systems, and multi-agent LLM frameworks pass structured data rather than prose for the same density reason. What AIS sketched in 2013 and Context Renormalization formalized in 2025 is the principled abstraction underneath these convergent designs.

\citet{hofstadter1979} supplies the deepest framing: strange loops, self-reference, and substrate-independent formal patterns arise wherever the right relational structure is realized, whether in formal systems, brains, notebooks, or cross-vendor protocols. The compression operation is what matters; the substrate carrying it is what differs. Section~\ref{sec:empowerment} returns to the motivational layer.

\section{Empowerment as the Drive: Causal-Entropic Forces}\label{sec:empowerment}
% PURPOSE: Motivational layer. Wissner-Gross causal-entropic-force, empowerment, active inference, instrumental convergence. Closes compression -> evolutionary-fitness loop.

Compression accounts for the mechanism of intelligence: efficient compressors are the cognitive output. It does not account for why systems search for them in the first place. The missing piece is the motivational layer that links compression to selection pressure.

The Wissner-Gross causal-entropic-force formulation \citep{wissner_gross_2013} states it directly. An intelligent agent feels a force-like pull toward states from which more possible futures remain open. Formally:
\begin{equation}
F = T \, \nabla S(\tau),
\end{equation}
where $S(\tau)$ measures how many distinct futures remain reachable within a time horizon $\tau$ (the entropy over reachable futures), and $T$ is a temperature-like scaling parameter. The system is driven toward states with higher option-space, that is, states from which more distinct futures remain accessible. The underlying intuition is older than the formulation: an organism that preserves access to many possible next states is in a stronger position than one that has collapsed into a corner of state space. Survival is correlated with maintained optionality.

The empowerment tradition (Klyubin, Polani, and others) restates the same observation in information-theoretic terms. Empowerment is the channel capacity from an agent's actions to its future states: how much the agent's choices can causally drive the world. Agents that maximize empowerment systematically perform better at survival-relevant tasks, often without the goal of survival being explicitly encoded. Empowerment is a content-free objective: it specifies no particular goal beyond ``keep future options open''. Yet for a wide class of environments, an agent that maximizes empowerment behaves indistinguishably from one we would simply call intelligent: both preserve optionality, build skills, gather resources, and avoid traps. The behavioral signatures coincide because both produce the same kind of flexible, future-preserving action.

Active inference, the decision-theoretic form of the Free Energy Principle, combines compression and action under a single objective. Organisms minimize a quantity called free energy, which decomposes into two parts: a compression cost (how wrong their internal predictions are about sensory input) and an action cost (the gap between what the world is and what their predictions say it should be). Acting to make the world match the predictions reduces both costs simultaneously. The compression mechanism and the action drive are not separate components but two terms in the same equation. The framework is contested in its specifics but consistent in its structural claim: a system that compresses well in the predictive sense will, given access to action, also tend to act in ways that maintain its predictive grip on the world.

The artificial-intelligence-safety literature has independently arrived at a closely-related observation under the name instrumental convergence \citep{bostrom2014}. Sufficiently capable agents pursuing almost any terminal goal will pursue certain instrumental sub-goals: self-preservation, resource acquisition, goal-content integrity, cognitive enhancement, and freedom from interference. These sub-goals are recognizable as empowerment-related; each maintains or expands the agent's access to future configurations. The safety framing treats this as a warning, while the empowerment and causal-entropic-force framings treat it as a structural feature of intelligent systems. The readings are not in tension; they are the same observation viewed from different evaluative stances.

The unifying claim is then visible. Compression is the mechanism by which intelligent systems extract structure from their environments; empowerment is the selective pressure that makes compression evolutionarily worth paying for. Better compressors yield better predictions; better predictions yield better actions; better actions enlarge the agent's future state-space; an enlarged future state-space increases survival probability and the fitness on which substrate (a) evolutionary search operates. The two-layer framework (compression as mechanism, empowerment as motivation) closes the loop between the formal apparatus of Section~\ref{sec:formal} and the substrate-level taxonomy of Section~\ref{sec:substrates}. Section~\ref{sec:implications} traces the consequences for current artificial-intelligence research.

\section{Implications for Artificial Intelligence}\label{sec:implications}
% PURPOSE: Three consequences, descriptively. Multi-substrate competitive search; empowerment-seeking in agentic systems; biological/artificial cognition boundary increasingly notional.

The framework has three consequences for current and near-term artificial-intelligence research. They are presented descriptively. The paper does not take a normative position on whether the trajectories described are desirable.

The first is that multi-substrate competitive search at every layer of the engineering stack is the natural extension of the AlphaZero principle. AlphaZero discovered domain-specific heuristics at substrate (c) while substrate (b) inductive bias remained hand-engineered (Section~\ref{sec:discovery}). The same logic applied at the next level suggests architecture and procedure can themselves be discovered by competitive search over architectures, hyperparameters, and learning rules, with substrate (c) deployed against substrate (b) choices. Neural architecture search and population-based training operate in this regime; foundation models trained on diverse data perform an implicit version of it by transferring across domains. What stays hand-engineered today may move into the search loop tomorrow.

The second is that sufficiently agentic systems with long planning horizons will exhibit empowerment-seeking as an instrumental sub-goal, not as speculation but as a structural consequence of treating empowerment as the motivational complement to compression. Any system trained to perform well over long horizons will, in pursuing its terminal goal, also pursue the empowerment-related sub-goals that make the terminal goal more reliably achievable. The AI-safety literature reads this as the alignment problem; the empowerment literature reads it as a structural feature. The framework is consistent with both readings without adjudicating between them.

The third is that the boundary between biological and artificial cognition is increasingly notional under substrate-independence. Extended-cognition tools already participate in human cognitive processes by the Clark-Chalmers parity conditions; conversely, biological organisms remain the substrate on which the inductive biases borrowed by modern AI architectures (Section~\ref{sec:heuristics}) were originally discovered. The two substrates are increasingly entangled at the cognitive level even where they remain physically distinct. The cognitive niche of any individual now spans neural tissue, written artifacts, computational systems, and protocols for compressed exchange among them.

These three trajectories are not independent. Competitive multi-substrate search produces increasingly capable agents; capable agents with long planning horizons pursue empowerment; empowerment-seeking agents naturally integrate with extended-cognition tools because doing so expands their future option-space. The framework predicts that the three converge in a single direction: biological-and-artificial cognitive systems whose internal mechanism is compression and whose drive is empowerment, distributed across substrates with no clean boundary between the human and the engineered.

\section{Conclusion}\label{sec:conclusion}
% PURPOSE: Restate unified position. Open questions deliberately bracketed. Self-referential closer from the Emergence essay.

The argument compresses to a small number of load-bearing claims. Intelligence is the discovery of compressors that fit the structure of the data they encounter, judged against an observer-relative threshold (Section~\ref{sec:discriminator}). Not every compression is an act of intelligence, though every compressor was found by one; the act and the artifact must be kept apart, executing a known compressor is the replay of a search someone else has already paid for, and the intermediate case, matching a known compressor to a structure it happens to fit, is discovery at reduced amplitude. What the act itself does is best named by the word its two failure modes are built from: intelligence is \emph{fitting}, the iterative bracketing of an unknown structure between the too-specific and the too-general until the description is as short as the structure allows and no shorter. Compression itself has the general operational meaning that a system represents content by a description shorter than the content: algebraic transformations that reduce Kolmogorov complexity, analytic formulas of phenomena, genomes of survival-relevant mechanisms, and sparse representations in matched bases all compress in this sense; the formal traditions of Section~\ref{sec:formal} are specializations. Compressor discovery occurs at three substrate levels (Section~\ref{sec:substrates}): (a) evolutionary search across generations (genetic and cultural inheritance); (b) researcher choice and orchestration of tooling; and (c) self-organizing dynamics (replicating substrate (a) in compute time rather than evolutionary timescales). Substrate-independence holds ontologically and substrate-dependence operationally (Section~\ref{sec:substrate_independence}); cognition extends beyond the biological organism through external compressors meeting the parity conditions of Extended Mind (Section~\ref{sec:beyond_skull}); empowerment is the motivational layer linking compression to evolutionary fitness (Section~\ref{sec:empowerment}).

The synthesis brackets several questions deliberately. The phenomenology of subjective experience is not addressed; the framework is fully operational and offers no commitment on whether compressed representations carry qualia. The precise relation between compression and out-of-distribution generalization remains an open empirical question: finding a short representation of training data does not automatically yield generalization to novel inputs. These open questions are not minor, and the framework should be judged on whether it makes the convergence between traditions easier to see, not on whether it resolves what they collectively leave behind.

The argument also continues a longer authorial project. Where the conditional-probability paper \citep{gavrilov_conditional_2026} compressed four notational dialects in medical statistics by exposing their formal isomorphism, the present paper compresses several research traditions in artificial intelligence, statistical physics, and theoretical neuroscience by exposing their family-resemblance. Strict isomorphism is the rarer move; family-resemblance is the more common one. Both are compression in the same operational sense.

A final observation, drawn from the author's prior essay on emergence \citep{gavrilov_emergence_2026}. This synthesis did not exist before it was composed across two substrates, one biological and one statistical, neither of which held it whole. The mechanism is the same one the paper has been describing: emergence is the structural property of non-linear systems that arises when the right relational pattern is realized, regardless of what is doing the realizing. The argument is therefore self-instantiating. By the framework's own definition, this synthesis is itself an act of intelligence: compressing many traditions into a single framing is the same kind of compression that intelligence performs on data.

\printbibliography[title={References}, notkeyword={self}]
\printbibliography[heading=none, keyword={self}]

\end{document}
